% called by main.tex
%
\chapter{Despliegue}
\label{ch::despliegue}

Este capítulo, opcional en la plantilla, documenta la \emph{estrategia} de despliegue. La
decisión más importante es deliberada y se explica a continuación: el sistema \textbf{no se
despliega} con capital real, sino que se detiene en una fase de seguimiento en sombra
(\emph{paper-shadow}).

\section{Estrategia de despliegue y seguimiento en paper-shadow}

Esta sección justifica la decisión de no desplegar el sistema en producción, describe la
infraestructura de seguimiento en sombra que se utiliza en su lugar y esboza qué requeriría
un despliegue futuro. La decisión de detenerse es coherente con la prudencia metodológica que
guía todo el trabajo.

\subsection{Por qué no se despliega}

Desplegar una política de decisión económica con cinco días de soporte fuera de muestra
sería metodológicamente imprudente, por mucho que el intervalo de confianza sea estrecho y
los cinco días positivos. Cinco días no constituyen una validación definitiva, y un sistema
que opera con capital real es especialmente sensible al sobreajuste cuando el horizonte de
datos es reducido. Por coherencia con la tesis del trabajo —que la honestidad del protocolo
prevalece sobre cualquier número concreto— el proyecto se detiene antes de operar.

\subsection{Seguimiento \emph{paper-shadow}}

En su lugar, el sistema se mantiene en \emph{paper-shadow}: se registra, operación a
operación, lo que la política \emph{habría} hecho, sin arriesgar capital. Esta
infraestructura produce una contabilidad detallada, un resumen diario y, sobre todo, una
puerta de decisión cuantitativa que determina cuándo —si alguna vez— tendría sentido pasar a
una fase superior. Las reglas de promoción son deliberadamente conservadoras y exigen
soporte temporal, no solo neto positivo:

\begin{itemize}[noitemsep]
  \item \textbf{Candidato a paper-shadow:} neto superior a 0,5 ticks, al menos 200
  operaciones y al menos 8 días de soporte.
  \item \textbf{Candidato a despliegue:} neto superior a 0,5 ticks, al menos 400 operaciones
  y al menos 12 días de soporte.
\end{itemize}

Con los datos actuales (318 operaciones, 5 días), el candidato no alcanza el primer umbral
por falta de días, y queda explícitamente en estado \emph{prometedor pero con datos
insuficientes}.

\subsection{Diseño de una hipotética API de inferencia}

A título de trabajo futuro, y solo si se superasen las puertas de promoción anteriores, un
despliegue eventual tomaría la forma de un servicio de inferencia con cuatro etapas
encadenadas. Primero, una etapa de ingesta que reconstruyese en tiempo real el estado del
mercado sobre la malla de dos segundos a partir de las mismas fuentes de captura. Segundo,
una etapa de cálculo que generase las 36 características causales del instante actual,
idéntica a la del entrenamiento para evitar discrepancias entre entrenamiento y producción
(\emph{train/serve skew}). Tercero, una etapa de decisión que consultase las dos cabezas del
especialista y devolviese una recomendación (operar o no) junto con el valor esperado y la
probabilidad de salud. Y cuarto, una capa de guardas que aplicase, antes de cualquier
recomendación, los filtros estrictos de entorno, el filtro de volatilidad y un mecanismo de
inhibición que cortase la operativa ante datos obsoletos o degradados.

Por encima de ese servicio operaría una capa de monitorización que vigilase la deriva de
distribución (especialmente el régimen de volatilidad), registrase cada decisión para
auditoría y dispusiese de un interruptor de parada inmediata. Conviene insistir en que esta
arquitectura es un \emph{diseño}, no una implementación: su desarrollo solo tendría sentido
tras acumular el soporte temporal que hoy falta, y siempre manteniendo el principio de
detenerse ante cualquier señal de degradación.
